\begin{blocksection}
\question \textbf{Session records.} What would Python display after executing
the following definitions? Read the field annotations before tracing.

\begin{lstlisting}
from dataclasses import dataclass

@dataclass
class Session:
    topic: str
    room: int
    minutes: int = 60

morning = Session("Recursion", 310)
afternoon = Session("Sequences", 205, 45)
\end{lstlisting}
\end{blocksection}

\begin{blocksection}
\begin{lstlisting}
>>> morning.topic
\end{lstlisting}
\begin{solution}[.35in]
\begin{lstlisting}
'Recursion'
\end{lstlisting}
\end{solution}
\end{blocksection}

\begin{blocksection}
\begin{lstlisting}
>>> morning.minutes
\end{lstlisting}
\begin{solution}[.35in]
\begin{lstlisting}
60
\end{lstlisting}
\end{solution}
\end{blocksection}

\begin{blocksection}
\begin{lstlisting}
>>> afternoon.room
\end{lstlisting}
\begin{solution}[.35in]
\begin{lstlisting}
205
\end{lstlisting}
\end{solution}
\end{blocksection}

\begin{blocksection}
\begin{lstlisting}
>>> afternoon.minutes
\end{lstlisting}
\begin{solution}[.35in]
\begin{lstlisting}
45
\end{lstlisting}
\end{solution}
\end{blocksection}

\begin{blocksection}
\begin{lstlisting}
>>> morning.minutes + afternoon.minutes
\end{lstlisting}
\begin{solution}[.35in]
\begin{lstlisting}
105
\end{lstlisting}
\end{solution}
\end{blocksection}

\begin{blocksection}
\begin{lstlisting}
>>> morning
\end{lstlisting}
\begin{solution}[.35in]
\begin{lstlisting}
Session(topic='Recursion', room=310, minutes=60)
\end{lstlisting}
\end{solution}
\end{blocksection}

\begin{blocksection}
\begin{lstlisting}
>>> Session(room=220, topic="Linked lists").topic
\end{lstlisting}
\begin{solution}[.35in]
\begin{lstlisting}
'Linked lists'
\end{lstlisting}
\end{solution}
\end{blocksection}

\begin{questionmeta}
\textbf{Teaching Tips}\par\nopagebreak[4]
Suggested Time: 6--8 min; Difficulty: Easy
\nopagebreak[4]
\begin{itemize}
    \item \textbf{Fields and types:} \lstinline{topic} describes text, while \lstinline{room} and \lstinline{minutes} describe integers. Annotations communicate expected types.
    \item \textbf{Construction:} \lstinline{@dataclass} supplies a constructor from the declared fields. Positional arguments follow field order; keyword arguments name the fields explicitly.
    \item \textbf{Defaults:} Omitting \lstinline{minutes} uses 60. Supplying 45 uses that value for the constructed session.
    \item \textbf{Access and display:} Dot notation reads one field. The generated representation names the class and its fields. Focus on constructing and reading instances; students do not need to implement special methods for this exercise.
\end{itemize}
\end{questionmeta}
